\question[5]
a. Schreibe eine Klasse \texttt{Node}, die einen Knoten eines binären Baums repräsentiert.

b. Schreibe eine Funktion \texttt{baum}, die den abgebildeten Baum zurückgibt.
Nutze dazu deine Node-Klasse. Verschachtele die Node-Aufrufe aus Lesbarkeitsgründen
höchstens einmal und gib den temporären Teilbäumen sprechende Namen.

\includegraphics[height=4cm]{\pfad/baumdef_01.png}

\begin{solutionbox}{6cm}
\begin{lstlisting}
class Node:
    def __init__(self, val, left=None, right=None):
        self.val = val
        self.left = left
        self.right = right

def baum():
    b10 = Node(10,Node(2),Node(3))
    b6 = Node(6,None,b10)
    b9 = Node(9,Node(1),None)
    b7 = Node(7,b6,b9)
    return b7
\end{lstlisting}
\end{solutionbox}
